% =====================================================================
\section{The Walkable Bunker}\label{sec:walk}
The command bunker is underground, so the player has to walk into it. Two pieces implement this: the \code{WalkSpace}
(floors, walls and indoor volumes, \code{scene/WalkSpace.h}) built by \code{Bunker.cpp}, and the \emph{walk camera}
(\code{Camera::moveWalk}). Both are CPU-only and unit-tested (T50--T54, T60--T69, T120--T126).

\subsection{WalkSpace}
\begin{itemize}
  \item \textbf{Floors.} A \code{WalkFloor} is a rectangle with a floor height at its low edge and one at its high edge along a slope axis (flat, along $x$ or along $z$) and a ceiling
  height. The bunker defines 12: the entrance ramp (a linear slope along $z$ from the surface at $z=13$ to the floor level $-3.4$\,m at $z=6$), a level
  approach, the door zone (with ceiling), the tunnel, the opening to the CIC, the command room ($x\in[-5.2,5.2]$, $z\in[-8.2,-2.1]$), the bunk room and its doorway, the generator room and its
  doorway, the shaft entrance and the exit shaft (a floor whose ceiling is the hatch height). Outside every floor the walker stands on the terrain.
  \item \textbf{Which floor.} \code{floorHeight(x,z,feetY)} returns the \emph{highest} floor region that contains $(x,z)$, lies at most a 0.45\,m step above the feet,
  and whose ceiling is still above the feet. A walker on the earth mound above the rooms therefore stays on the terrain, and a walker in the tunnel is not pulled up onto the roof.
  \item \textbf{Walls.} 36 axis-aligned wall boxes (retaining walls, cross walls, consoles and other furniture, door jambs). \code{collide()} pushes a vertical cylinder (radius 0.3\,m, feet to
  head) out of every wall box and every extra box, in two passes to settle corners; boxes that end below the feet or begin above the head are ignored.
  \item \textbf{Dynamic obstacles.} The closed blast door and a lowered barrier arm are passed as extra boxes every frame (T64, T121).
  \item \textbf{Indoor volumes.} Five boxes (tunnel, CIC, bunk, generator, shaft) are uploaded to the shader as the indoor mask (Section~\ref{sec:shading}); the renderer reads the first eight.
\end{itemize}
The bunker part budget in test T66 reads: \bunkerBudget. T69: 12 floors, 36 walls, 5 indoor volumes.

\subsection{Walk camera}
Key \texttt{C} cycles Orbit, Free, Chase, UAV CAM and Walk (T55). In walk mode the eye is 1.7\,m above the walker's feet (T51: maximum deviation 0.013\,m over 2\,400 steps), the field of
view is $60^\circ$, the speed is 1.6\,m/s (4.5\,m/s with Shift) and the body radius 0.3\,m. Long steps are split into $\le0.15$\,m pieces so a fast frame cannot tunnel through a door; after each piece
the feet follow the floor (down the ramp: T120, maximum feet deviation 0\,m). The eye is limited to 0.15\,m below the ceiling (T123); a mouse-look pitch is clamped to $\pm80^\circ$
(T54); the position is clamped to a 52\,m radius (T53) and the walker cannot fall into the open trench from the side (T52). A small head bob ($0.025\,\text{m}\cdot\sin$ of a speed-driven
phase) is added to the rendered eye only.

\subsection{Verified walk-through}
T62 walks from the ramp top to the CIC centre in 367 steps with 0 wall contacts, a maximum step of 0.023\,m and a minimum head room of 2.60\,m. T121: with the doors closed the walker stops at
$z=4.4$\,m (the door plane is at $z=4$) and never leaves the ramp ($|x|\le1.2$\,m). T122: with the doors open he reaches the CIC ($z=-6.0$\,m, feet $-3.40$\,m, smallest ceiling gap 0.90\,m).
T126: the director's bunker shot descends the ramp, passes the door plane and ends in the rooms with a maximum vertical step of 0.118\,m, and T129 checks that the
director opens the doors early. Walking up to a door and pressing \texttt{Enter} within 3.5\,m opens or closes it, and the door status appears in the HUD.

\begin{figure}[h]
  \centering
  \begin{subfigure}{0.49\linewidth}\picfour{bunker_ramp}{\linewidth}\caption{Entrance ramp (walk camera)}\end{subfigure}
  \begin{subfigure}{0.49\linewidth}\picfour{bunker_door}{\linewidth}\caption{Blast doors}\end{subfigure}\\[1mm]
  \begin{subfigure}{0.49\linewidth}\picfour{bunker_tunnel}{\linewidth}\caption{Tunnel}\end{subfigure}
  \begin{subfigure}{0.49\linewidth}\picfour{bunker_cic}{\linewidth}\caption{Command room (CIC)}\end{subfigure}
  \caption{Walking into the bunker (captures \code{bunker\_*}; eye heights 1.7\,m).}
  \label{fig:bunkerwalk}
\end{figure}

\begin{figure}[h]
  \centering
  \begin{subfigure}{0.49\linewidth}\picthree{walk_bunker}{\linewidth}\caption{v3 walk camera at the old bunker}\end{subfigure}
  \begin{subfigure}{0.49\linewidth}\picthree{walk_forest}{\linewidth}\caption{v3 walk camera in the forest}\end{subfigure}
  \caption{The walk camera introduced in v3 (EPIC captures).}
\end{figure}
